% Generated draft: requires mathematical and source review.
\documentclass[11pt]{article}
\usepackage[margin=25mm]{geometry}
\usepackage{fontspec}
\setmainfont{Latin Modern Roman}
\usepackage{amsmath,amssymb,mathtools,unicode-math}
\setmathfont{Latin Modern Math}
\usepackage{enumitem,needspace,xurl,hyperref}
\setlength{\parindent}{0pt}
\setlength{\parskip}{5pt}
\setlength{\emergencystretch}{3em}
\newcommand{\N}{\mathbb{N}}
\newcommand{\Z}{\mathbb{Z}}
\newcommand{\Q}{\mathbb{Q}}
\newcommand{\R}{\mathbb{R}}
\newcommand{\C}{\mathbb{C}}
\newcommand{\F}{\mathbb{F}}
\newcommand{\A}{\mathbb{A}}
\newcommand{\PP}{\mathbb{P}}
\newcommand{\E}{\mathbb{E}}
\newcommand{\eps}{\varepsilon}
\newcommand{\dd}{\,\mathrm{d}}
\begin{document}
\section*{UnsolvedMath ID: 11000032}
\textbf{Source problem code:} catalogue entry
\subsection{Definitions and mathematical statement}
Fix the source's genus \(g\), and let \(\operatorname{Mod}_g\) be the corresponding mapping class group. The question is whether the isomorphism problem for its \emph{finitely presented subgroups} is solvable: given effective descriptions of two such subgroups \(H,K\leq\operatorname{Mod}_g\), decide whether \(H\cong K\) as abstract groups. The decision is not whether they are conjugate inside \(\operatorname{Mod}_g\), and no embedding-preserving isomorphism is required.

One precise presentation-input version uses
\[\mathcal P_g=\{P:\ P\text{ is a finite group presentation and }G(P)\text{ embeds in }\operatorname{Mod}_g\}.\]
It asks for a single terminating algorithm on inputs \(P,Q\in\mathcal P_g\) returning yes exactly when \(G(P)\cong G(Q)\), under the promise that both inputs lie in this class. Another possible version gives finite tuples of mapping-class words generating \(H,K\), with a promise that the generated subgroups are finitely presented, or also supplies presentations and embedding maps. The source does not specify which input model is intended; it cannot be assumed that these data can be converted effectively with no additional work. Behavior on inputs outside the promised class need not be prescribed unless a total-recognition problem is intended.

The source contrasts this with reported unsolvability for all finitely generated linear groups, and even for finitely generated subgroups of \(\operatorname{GL}(n,\Z)\). Those are stated for finitely generated classes, not the narrower finitely presented class in the main question; the exact \(n\), encodings and supporting theorems are not provided. No transfer of those negative results to the requested mapping-class subgroup problem is asserted. The ending introduces another algorithmic question but does not supply it.
\subsection{Short English statement}
Is there an algorithm deciding abstract isomorphism for finitely presented subgroups of a mapping class group? The subgroup encoding and fixed-versus-variable genus must be specified; the source's negative linear-group comparisons concern broader finitely generated classes.
\subsection{Sources}
\begin{itemize}
\item Local UnsolvedMath registry, permanent ID \(11000032\), exact source code \texttt{catalogue entry}.
\item Farb, \emph{Problems on Mapping Class Groups and Related Topics}, Question 3.11, PDF page 34: \url{https://www.math.uchicago.edu/~farb/papers/mcgbook.pdf}. Registry access 2026-07-29; pdf-font-structure extraction.
\item The source states the linear-group unsolvability comparisons without a precise citation in this excerpt. The research summary supplies no resolving reference.
\end{itemize}
\end{document}
